\textbf{Phần 1. Trắc nghiệm lựa chọn}
\setcounterpageref{ex}{0}

\begin{ex} %[2D2N13]
Lan có 8 cái kẹo, Mai có nhiều hơn Lan 5 cái kẹo. Hỏi Mai có bao nhiêu cái kẹo?
\choice
{\True $13$ cái}
{$12$ cái}
{$14$ cái}
{$15$ cái}
\loigiai{
$8 + 5 = 13$ cái kẹo.\\
\textbf{Đáp án của câu hỏi là A}
}
\end{ex}

\begin{ex} %[2D2N13]
An có 15 viên bi, Bình có ít hơn An 6 viên bi. Hỏi Bình có bao nhiêu viên bi?
\choice
{$8$ viên}
{$10$ viên}
{$11$ viên}
{\True $9$ viên}
\loigiai{
$15 - 6 = 9$ viên bi.\\
\textbf{Đáp án của câu hỏi là D}
}
\end{ex}

\begin{ex} %[2D2H13]
Bể thứ nhất có 9 con cá, bể thứ hai có nhiều hơn bể thứ nhất 7 con cá. Bể thứ hai có số con cá là
\choice
{\True $16$ con}
{$15$ con}
{$17$ con}
{$18$ con}
\loigiai{
$9 + 7 = 16$ con cá.\\
\textbf{Đáp án của câu hỏi là A}
}
\end{ex}

\begin{ex} %[2D2H13]
Tổ 1 có 14 bạn, tổ 2 có ít hơn tổ 1 là 5 bạn. Tổ 2 có số bạn là
\choice
{$8$ bạn}
{$10$ bạn}
{$11$ bạn}
{\True $9$ bạn}
\loigiai{
$14 - 5 = 9$ bạn.\\
\textbf{Đáp án của câu hỏi là D}
}
\end{ex}

\begin{ex} %[2D2H13]
Bác Ba nuôi 8 con bò, số con lợn nhiều hơn số con bò là 9 con. Bác Ba nuôi số con lợn là
\choice
{$16$ con}
{\True $17$ con}
{$18$ con}
{$15$ con}
\loigiai{
$8 + 9 = 17$ con lợn.\\
\textbf{Đáp án của câu hỏi là B}
}
\end{ex}

\begin{ex} %[2D2H13]
Cây ổi có 16 quả, cây xoài có ít hơn cây ổi 7 quả. Cây xoài có số quả là
\choice
{$8$ quả}
{$10$ quả}
{$11$ quả}
{\True $9$ quả}
\loigiai{
$16 - 7 = 9$ quả.\\
\textbf{Đáp án của câu hỏi là D}
}
\end{ex}

\begin{ex} %[2D2V13]
Đoạn thẳng $AB$ dài $9\text{ cm}$, đoạn thẳng $CD$ dài hơn đoạn thẳng $AB$ là $6\text{ cm}$. Độ dài đoạn thẳng $CD$ là
\choice
{$14\text{ cm}$}
{$16\text{ cm}$}
{\True $15\text{ cm}$}
{$17\text{ cm}$}
\loigiai{
$9 + 6 = 15\text{ cm}$.\\
\textbf{Đáp án của câu hỏi là C}
}
\end{ex}

\begin{ex} %[2D2V13]
Sợi dây màu xanh dài $18\text{ cm}$, sợi dây màu đỏ ngắn hơn sợi dây màu xanh $9\text{ cm}$. Độ dài sợi dây màu đỏ là
\choice
{$8\text{ cm}$}
{\True $9\text{ cm}$}
{$10\text{ cm}$}
{$11\text{ cm}$}
\loigiai{
$18 - 9 = 9\text{ cm}$.\\
\textbf{Đáp án của câu hỏi là B}
}
\end{ex}

\begin{ex} %[2D2V13]
Hùng có 9 viên bi, Dũng có nhiều hơn Hùng 8 viên bi. Cả hai bạn có tất cả bao nhiêu viên bi?
\choice
{\True $26$ viên}
{$25$ viên}
{$27$ viên}
{$17$ viên}
\loigiai{
Dũng có: $9 + 8 = 17$ viên.\\
Cả hai bạn có: $9 + 17 = 26$ viên bi.\\
\textbf{Đáp án của câu hỏi là A}
}
\end{ex}

\begin{ex} %[2D2V13]
Hoa có 14 bông hoa, Cúc có ít hơn Hoa 6 bông hoa. Cả hai bạn có tất cả bao nhiêu bông hoa?
\choice
{$21$ bông}
{$23$ bông}
{\True $22$ bông}
{$24$ bông}
\loigiai{
Cúc có: $14 - 6 = 8$ bông.\\
Cả hai bạn có: $14 + 8 = 22$ bông hoa.\\
\textbf{Đáp án của câu hỏi là C}
}
\end{ex}

\begin{ex} %[2D2V13]
Nam có 8 cái bút chì, số bút chì của Minh nhiều hơn Nam là 4 cái. Minh có số bút chì là
\choice
{$11$ cái}
{\True $12$ cái}
{$13$ cái}
{$14$ cái}
\loigiai{
$8 + 4 = 12$ cái.\\
\textbf{Đáp án của câu hỏi là B}
}
\end{ex}

\begin{ex} %[2D2C13]
Có hai ngăn sách. Ngăn thứ nhất có 9 quyển sách. Nếu chuyển 3 quyển từ ngăn thứ hai sang ngăn thứ nhất thì hai ngăn bằng nhau. Hỏi ban đầu ngăn thứ hai có bao nhiêu quyển sách?
\choice
{$14$ quyển}
{$16$ quyển}
{\True $15$ quyển}
{$12$ quyển}
\loigiai{
Sau khi chuyển, ngăn 1 có: $9 + 3 = 12$ quyển.\\
Vậy lúc sau ngăn 2 cũng có 12 quyển $\Rightarrow$ Ban đầu ngăn 2 có: $12 + 3 = 15$ quyển.\\
\textbf{Đáp án của câu hỏi là C}
}
\end{ex}

\textbf{Phần 2. Trắc nghiệm đúng sai}
\setcounterpageref{ex}{0}

\begin{ex} %[2D2H13]
Tổ 1 có 8 bạn, tổ 2 có nhiều hơn tổ 1 là 5 bạn. Xét tính đúng sai:
\choiceTF
{\True Tìm số bạn tổ 2 ta làm phép tính cộng ($8 + 5$).}
{\True Tổ 2 có 13 bạn.}
{Cả hai tổ có tất cả 20 bạn.}
{\True Cả hai tổ có tất cả 21 bạn ($8 + 13 = 21$).}
\loigiai{
\begin{itemize}
    \item Phép tính cộng $8 + 5$: Đúng.
    \item Tổ 2 có 13 bạn: Đúng.
    \item Tổng: $8 + 13 = 21 \neq 20$. Ý 3 Sai, Ý 4 Đúng.
\end{itemize}
\begin{itemize}
    \item a) \textbf{Đúng}.
    \item b) \textbf{Đúng}.
    \item c) \textbf{Sai}.
    \item d) \textbf{Đúng}.
\end{itemize}
}
\end{ex}

\begin{ex} %[2D2V13]
Thùng to có 15 lít dầu, thùng nhỏ có ít hơn thùng to 7 lít dầu. Xét tính đúng sai:
\choiceTF
{\True Thùng nhỏ có 8 lít dầu ($15 - 7 = 8$).}
{\True Cả hai thùng có $23\text{ lít}$ dầu ($15 + 8 = 23$).}
{Thùng nhỏ có nhiều hơn 10 lít dầu.}
{\True Thùng nhỏ có ít hơn 10 lít dầu (chỉ có $8\text{ lít}$).}
\loigiai{
\begin{itemize}
    \item Thùng nhỏ: $15 - 7 = 8\text{ l}$: Đúng.
    \item Tổng: $15 + 8 = 23\text{ l}$: Đúng.
    \item $8 < 10$ (không nhiều hơn). Ý 3 Sai, Ý 4 Đúng.
\end{itemize}
\begin{itemize}
    \item a) \textbf{Đúng}.
    \item b) \textbf{Đúng}.
    \item c) \textbf{Sai}.
    \item d) \textbf{Đúng}.
\end{itemize}
}
\end{ex}

\textbf{Phần 3. Trả lời ngắn}
\setcounterpageref{ex}{0}

\begin{ex} %[2D2H13]
Lan có 9 cái kẹo, Huệ có nhiều hơn Lan 6 cái kẹo. Hỏi Huệ có bao nhiêu cái kẹo?
\shortans{15}
\loigiai{
\textbf{Đáp số:} $15$\\
$9 + 6 = 15$ cái.
}
\end{ex}

\begin{ex} %[2D2H13]
Đàn vịt có 17 con, đàn gà có ít hơn đàn vịt 9 con. Hỏi đàn gà có bao nhiêu con?
\shortans{8}
\loigiai{
\textbf{Đáp số:} $8$\\
$17 - 9 = 8$ con.
}
\end{ex}

\begin{ex} %[2D2V13]
Đoạn dây thứ nhất dài $8\text{ dm}$, đoạn dây thứ hai dài hơn đoạn thứ nhất $7\text{ dm}$. Tính độ dài đoạn dây thứ hai theo đề-xi-mét.
\shortans{15}
\loigiai{
\textbf{Đáp số:} $15$\\
$8 + 7 = 15\text{ dm}$.
}
\end{ex}

\begin{ex} %[2D2V13]
Hộp màu xanh có 14 cái bút, hộp màu đỏ có ít hơn hộp màu xanh 8 cái bút. Hỏi hộp màu đỏ có bao nhiêu cái bút?
\shortans{6}
\loigiai{
\textbf{Đáp số:} $6$\\
$14 - 8 = 6$ cái.
}
\end{ex}

\begin{ex} %[2D2V13]
Bác An nuôi 9 con thỏ trắng, số thỏ xám nhiều hơn thỏ trắng 8 con. Hỏi bác An nuôi bao nhiêu con thỏ xám?
\shortans{17}
\loigiai{
\textbf{Đáp số:} $17$\\
$9 + 8 = 17$ con.
}
\end{ex}

\begin{ex} %[2D2C13]
Mai có 8 bông hoa. Nếu Mai hái thêm 4 bông hoa nữa thì Mai vẫn ít hơn Lan 3 bông hoa. Hỏi Lan có bao nhiêu bông hoa?
\shortans{15}
\loigiai{
\textbf{Đáp số:} $15$\\
$8 + 4 + 3 = 15$ bông hoa.
}
\end{ex}
